\section{Integrierte Schlussfolgerung v1.6: drei innere Richtungen sind kein Dimensionsbeweis}
Die neue Eingabe verbindet $\mu_4$, eine zweidimensionale komplexe Darstellung, drei Pauli-Richtungen und eine räumliche Dreidimensionalität. Diese Kette enthält zusätzliche Voraussetzungen: Die komplexen irreduziblen Darstellungen von $\mu_4$ selbst sind eindimensional. Für $\C^2$ braucht man weitere nichtkommutierende Operationen. Drei Pauli-Koeffizienten begrenzen nicht automatisch die Dimension ihres Ortsarguments. Ein regulärer isolierter Weyl-Punkt in einer zweibandigen Beschreibung ist ein präziser bedingter Test, keine Auswahl der ursprünglichen Raumdimension.

Die $A_3$-Fundamentalgewichte bilden tatsächlich ein Tetraeder in einem dreidimensionalen Gewichtsraum; seine Gewichts- und Wurzelgitter haben BCC- beziehungsweise FCC-Lesarten. Die neuen Prüfungen trennen diese exakte Geometrie von der Wahl physischer Nachbarschaften. Ein ausgewählter Lift in das volle $E_8$-Gitter muss zusätzlich Metrik, Grad, geschlossene Wege und native Operatorprodukte erhalten. Die hierfür gefundenen Hindernisse und der inzwischen positiv berechnete interne Walk-Baustein stehen in Kapitel~\ref{ch:v16-operations}.
